% VIGTIGT! For at dokumentet kan compile, skal man benytte XeLaTeX eller LuaLaTeX som compiler. Dette kan gøres i Overleaf ved at gå til Menu -> Settings -> Compiler -> Vælg XeLaTeX/LuaLaTeX
\documentclass[t, 10pt, aspectratio=169]{beamer}

\usepackage{KUstyle}
\usepackage{subfiles}
\usepackage{listings}

\toplinje{Bachelor Defence} %Teksten på toplinjen. Hvis ingen tekst ønskes, skal kommandoen fjernes

\begin{document}

% Det første slide. Her kan man f.eks. ændre hovedtitel, undertitlen, oplægsholder, KU-enhed og dato
{
\setbeamertemplate{background}{\includegraphics[width=\paperwidth,height=\paperheight]{KU/forside.pdf}}
\begin{frame}
	\begin{textblock*}{\textwidth}(0\textwidth,0.1\textheight)
		\begin{beamercolorbox}[wd=7.8cm,ht=7.3cm,sep=0.5cm]{hvidbox}
			\fontsize{5}{10}\fontfamily{lmr}\selectfont \textls[200]{KØBENHAVNS UNIVERSITET}
			\noindent\textcolor{KUrod}{\rule{6.8cm}{0.4pt}}
		\end{beamercolorbox}
	\end{textblock*}
	\begin{textblock*}{\textwidth}(0\textwidth,0.1\textheight)
		\begin{beamercolorbox}[wd=7.8cm,sep=0.5cm]{hvidbox}
			\Huge \textcolor{KUrod}{RISC-V based computers in the data
				center}
			\vspace{0.5cm}
			\par
			\Large Bachelor Defence
			\vspace{0.5cm}
			\par
			\normalsize Simon Vinding Brodersen, 20-06-2024
		\end{beamercolorbox}
	\end{textblock*}
	\begin{textblock}{1}(6.3,11.38)
		\includegraphics[width=1cm]{KU/KU-logo.png}
	\end{textblock}
\end{frame}
}

\subfile{./sections/Introduction.tex}
\subfile{./sections/Background.tex}
\subfile{./sections/Design.tex}
\subfile{./sections/Implementation.tex}
\subfile{./sections/Evaluation.tex}
\subfile{./sections/RelatedWork.tex}
\subfile{./sections/Conclusion.tex}

\end{document}
